\chapter{Giới Thiệu}
\label{chap:intro}

% ─────────────────────────────────────────────────────────────────────────────
\section{Đặt Vấn Đề}
\label{sec:motivation}

Thị trường tuyển dụng nhân lực công nghệ thông tin tại Việt Nam và trên thế giới đang chứng kiến sự cạnh tranh ngày càng gay gắt. Theo khảo sát của TopDev (2024) \cite{topdev2024}, hơn 70\% ứng viên IT tại Việt Nam thừa nhận chưa chuẩn bị đủ kỹ năng kỹ thuật cần thiết trước các vòng phỏng vấn, trong khi phần lớn các nền tảng luyện phỏng vấn hiện có (LeetCode, HackerRank, Pramp) chủ yếu tập trung vào lập trình thuật toán, thiếu nội dung dành riêng cho ba vai trò dữ liệu phổ biến nhất: Data Analyst (DA), Data Scientist (DS) và Data Engineer (DE).

Hai thách thức cốt lõi tồn tại với các nền tảng luyện phỏng vấn hiện tại:

\textbf{(1) Thiếu cá nhân hóa theo năng lực.} Hầu hết hệ thống cung cấp câu hỏi theo cách tĩnh (static), không phản ánh được trình độ thực tế của từng ứng viên. Người dùng có kinh nghiệm phải làm lại các câu hỏi quá dễ, trong khi người mới bắt đầu có thể nản lòng với câu hỏi ngoài tầm. Vấn đề này tương ứng với khái niệm \textit{zone of proximal development} trong tâm lý học giáo dục \cite{vygotsky1978mind} -- câu hỏi hiệu quả nhất là câu hỏi hơi khó hơn năng lực hiện tại của người học một chút.

\textbf{(2) Không theo dõi được quỹ đạo kiến thức.} Hệ thống hiện có không duy trì và cập nhật trạng thái kiến thức (knowledge state) của người học theo thời gian. Do đó, sau mỗi phiên luyện tập, hệ thống không biết người dùng đã tiến bộ ở điểm nào, còn yếu ở kỹ năng nào, và câu hỏi tiếp theo nên theo hướng nào.

Knowledge Tracing (KT) \cite{corbett1994knowledge} là hướng giải quyết được nghiên cứu nhiều trong lĩnh vực Educational Data Mining, giúp mô hình hóa và dự đoán trạng thái kiến thức của người học qua lịch sử tương tác. Các mô hình KT từ Bayesian Knowledge Tracing \cite{corbett1994knowledge} đến Deep Knowledge Tracing \cite{piech2015deep} và Self-Attentive Knowledge Tracing \cite{pandey2019self} đã chứng minh hiệu quả trong môi trường học thuật (MOOC, bài tập toán, ngôn ngữ). Tuy nhiên, việc áp dụng KT cho bài toán luyện phỏng vấn IT -- với đặc thù câu hỏi mở, chấm điểm tự động bằng LLM và nhiều loại câu hỏi (lý thuyết, coding, tình huống) -- vẫn là một bài toán mở.

Luận văn này đề xuất và xây dựng \system{}, một hệ thống luyện phỏng vấn IT tích hợp KT, nhằm giải quyết đồng thời hai thách thức trên: theo dõi năng lực tự động và gợi ý câu hỏi cá nhân hóa theo từng vai trò nghề nghiệp.

% ─────────────────────────────────────────────────────────────────────────────
\section{Mục Tiêu Nghiên Cứu}
\label{sec:objectives}

Luận văn đặt ra bốn mục tiêu chính:

\begin{enumerate}
  \item \textbf{Xây dựng bộ dữ liệu luyện KT phù hợp cho bài toán phỏng vấn IT.} Do thiếu dữ liệu thực, luận văn thiết kế pipeline mô phỏng dữ liệu tổng hợp có tính chất thực tế cao, kết hợp IRT \cite{embretson2000item} để sinh hành vi người dùng và LLM để xây dựng question bank đa dạng cho ba vai trò DA/DS/DE.

  \item \textbf{Huấn luyện và so sánh các mô hình KT.} Triển khai và đánh giá ba họ mô hình -- BKT \cite{corbett1994knowledge}, DKT \cite{piech2015deep}, SAKT \cite{pandey2019self} -- với các biến thể sử dụng nhãn nhị phân và quality label liên tục. Phân tích điểm mạnh/yếu của từng mô hình trên tập dữ liệu phỏng vấn IT.

  \item \textbf{Xây dựng engine gợi ý dựa trên KT.} Thiết kế cơ chế cập nhật skill vector bằng cách kết hợp KT prediction (70\%) và EMA (30\%), tích hợp vào pipeline đề xuất câu hỏi theo chiến lược zone of proximal development, có hỗ trợ Computerized Adaptive Testing (CAT) hai giai đoạn \cite{wainer2000computerized}.

  \item \textbf{Tích hợp vào ứng dụng web hoàn chỉnh.} Xây dựng \system{} -- giao diện React SPA được phục vụ qua Streamlit, kết nối với FastAPI backend -- hỗ trợ chấm điểm tự động bằng Gemini LLM và fallback local rubric \cite{team2023gemini}.
\end{enumerate}

% ─────────────────────────────────────────────────────────────────────────────
\section{Phạm Vi Nghiên Cứu}
\label{sec:scope}

\textbf{Trong phạm vi:}
\begin{itemize}
  \item Ba vai trò nghề nghiệp: Data Analyst (DA), Data Scientist (DS), Data Engineer (DE).
  \item 17 Knowledge Components (KC) phân bổ theo vai trò: DA (5 KC), DS (7 KC), DE (5 KC).
  \item Bộ câu hỏi gồm 5 loại: trắc nghiệm (MC\_SINGLE), đúng/sai (TRUE\_FALSE), điền chỗ trống (FILL\_BLANK), lý thuyết mở (THEORY), và coding (CODING/CODING\_EXERCISE).
  \item Dữ liệu training KT: 1{,}000 người dùng mô phỏng, $\sim$40{,}000 tương tác.
  \item Chấm điểm câu hỏi mở bằng Gemini 3.5 Flash với local rubric fallback.
\end{itemize}

\textbf{Ngoài phạm vi:}
\begin{itemize}
  \item Thu thập và sử dụng dữ liệu người dùng thực (do giới hạn thời gian và quyền riêng tư).
  \item Triển khai production với load balancing và autoscaling.
  \item Đánh giá A/B testing so với các hệ thống thương mại.
  \item Hỗ trợ các ngôn ngữ lập trình thực thi trong sandbox (chỉ hiển thị).
\end{itemize}

% ─────────────────────────────────────────────────────────────────────────────
\section{Đóng Góp Chính}
\label{sec:contributions}

Luận văn có bốn đóng góp chính:

\begin{enumerate}
  \item \textbf{Pipeline xây dựng dữ liệu phỏng vấn IT tổng hợp}: Một pipeline đầy đủ từ sinh câu hỏi (LLM), tạo người dùng ảo (correlated Gaussian skill), đến mô phỏng tương tác (IRT + power-based graded response). Pipeline này có thể tái sử dụng cho các bài toán KT trong lĩnh vực nghề nghiệp khác.

  \item \textbf{So sánh hệ thống KT với quality label}: Đề xuất và đánh giá biến thể DKT-Quality sử dụng điểm chất lượng liên tục $\in \{0, 0.5, 1\}$ thay vì nhãn nhị phân, cho thấy cải thiện đáng kể về RMSE và PR-AUC trên bộ dữ liệu phỏng vấn IT.

  \item \textbf{Cơ chế tích hợp KT--EMA cho real-time skill update}: Công thức blend KT(70\%) + EMA(30\%) với ngưỡng kích hoạt tối thiểu 3 tương tác, giải quyết bài toán cold-start và đảm bảo tính ổn định trong giai đoạn đầu.

  \item \textbf{Hệ thống luyện phỏng vấn IT end-to-end}: Ứng dụng \system{} hoàn chỉnh tích hợp KT, 2-Stage Branching CAT, chấm điểm LLM, và giao diện React với FastAPI backend -- tất cả trong một codebase duy nhất có thể chạy local.
\end{enumerate}

% ─────────────────────────────────────────────────────────────────────────────
\section{Cấu Trúc Luận Văn}
\label{sec:structure}

Phần còn lại của luận văn được tổ chức như sau:

\textbf{Chương \ref{chap:theory} -- Cơ Sở Lý Thuyết:} Trình bày nền tảng lý thuyết của Knowledge Tracing (BKT, DKT, SAKT), Item Response Theory, và hệ thống gợi ý trong ngữ cảnh luyện phỏng vấn. Đồng thời khảo sát các công trình liên quan trong lĩnh vực.

\textbf{Chương \ref{chap:data} -- Xây Dựng Dữ Liệu và Phân Tích Khám Phá:} Mô tả chi tiết pipeline xây dựng question bank và dữ liệu mô phỏng, kết hợp với phân tích EDA toàn diện về đặc tính bộ dữ liệu.

\textbf{Chương \ref{chap:experiment} -- Thực Nghiệm và Đánh Giá:} Trình bày thiết kế thực nghiệm, kết quả huấn luyện và đánh giá các mô hình KT, phân tích so sánh, và thảo luận về hạn chế.

\textbf{Chương \ref{chap:conclusion} -- Kết Luận và Hướng Phát Triển:} Tổng kết kết quả đạt được, đóng góp, và đề xuất hướng phát triển cho nghiên cứu tương lai.
